\documentclass[a4paper,spanish]{article}
\usepackage{lmodern}
\renewcommand{\familydefault}{\sfdefault} %% Letra cuerpo 
\usepackage[T1]{fontenc}
\usepackage[utf8]{inputenc}
\usepackage[spanish]{babel}
\usepackage{array}
\usepackage{float}
\usepackage{fancybox}
\usepackage{calc}
\usepackage{units}
\usepackage{textcomp}
\usepackage{multirow}
\usepackage{amsmath}
\usepackage{amssymb}
\usepackage{graphicx}
\usepackage{nomencl}
\usepackage{fancyhdr}
\usepackage{graphics}
\usepackage{circuitikz}
\usepackage{graphicx}
\usepackage{pgf,tikz}
\usepackage{multicol,caption}
\usepackage{mathtools}
\usepackage[makeroom]{cancel}
\usepackage{lscape}
\usepackage{rotating}
\usepackage{lastpage}
\usepackage{titlesec}
\usepackage{tcolorbox}
\usepackage[margin=1in,headsep=1cm,headheight= 1in,top=1.5in]{geometry}
\usetikzlibrary{arrows}
\usetikzlibrary{shapes.symbols}
\usetikzlibrary{shapes.arrows}
\usetikzlibrary{calc,intersections,through,backgrounds}
\usetikzlibrary{patterns}
\PassOptionsToPackage{caption=false}{subfig}
\usepackage{subfig}
\makeatother
\graphicspath{ {imagenes/} }
\newtheorem{ej}{{Ejemplo }}[section]
\pagestyle{fancy}

\titleformat{name=\section}[block]
{\sffamily\bfseries\Large}
{}
{0pt}
{\colorsection}
\titlespacing{\section}{0pt}{\baselineskip}{\baselineskip}

\newcommand{\colorsection}[1]{%
	\colorbox{teal!40}{\parbox{\dimexpr\textwidth-2\fboxsep}{\ #1}}}

\setcounter{secnumdepth}{0} %% sacar numeracion a las secciones
\setcounter{tocdepth}{4}


\begin{document}




\begin{titlepage}
	\centering
	{\includegraphics[width=0.7\textwidth]{logo}\par}
	\vspace{1cm}
	{\bfseries\LARGE Escuela de Educación Secundaria Técnica N\(^{ro}\) de Vicente López \\ Eduardo Ader \par}
	\vspace{1cm}
	{\Large Montaje de Proyectos Electrónicos - \(6^{to}\)Año \par}
	\vspace{3cm}
	{\Huge Proyecto Integrador \par}
	\vspace{2cm}

	{\Large Docentes: \par}
	{\Large Alejandro Hsia, Esteban Lemos, Carlos Robello, Guillermo Ruisi \par}
	\vfill

\end{titlepage}


\clearpage
%Indice automático
\fancyhead{}
\fancyhead[C]{ÍNDICE}

\tableofcontents
\clearpage
%formato del resto del documento


\rhead{\leftmark \\ \rightmark}

\renewcommand{\footrulewidth}{2pt}
\renewcommand{\headrulewidth}{2pt}


\rfoot{\thepage\ de \pageref{LastPage}}
\fancyhead{}
\cfoot{\footnotesize\textbf{ Escuela de Educación Secundaria Técnica N$^\circ$ 1 - Vicente López}\\Cerrito N$^\circ$ 3966 (1605) Villa Adelina, Partido de Vicente López.\\Tel:(011) 47350174 -- email: tecnica1vl@gmail.com}

\lhead{\begin{tikzpicture}
		\node (0,0)[]{\includegraphics[width=0.5in]{logo}};
		\node [text width=8cm,align=left,xshift=5cm] at (0,0){\textbf{E.E.S.T. N$^\circ$ 1 V.L.}\linebreak MPE\linebreak PROYECTO INTEGRADOR} ;
	\end{tikzpicture}
}
\rhead{\leftmark \\\rightmark}


\newpage



\part{Proyecto Integrador}

\section{Objetivos del trabajo}

En este trabajo práctico integrador, se les desafía a diseñar, implementar y probar un sistema de control de temperatura que utiliza dos cargas eléctricas distintas para realizar dos funciones opuestas: enfriamiento y calefacción. El objetivo principal es lograr un control preciso de la temperatura de un entorno dado.

\section{Desarrollo}
\begin{enumerate}
	\item \textit{Diseño del Sistema:} Deben diseñar un sistema de control mediante una placa, que encienda dos cargas de 220 V de CA separadas, una para enfriar (ventilador con motor monofasico) y otra para calentar (resistencia de 2kW).
Las mismas serán activadas por dos niveles lógicos (1 o 0) que provienen del microcontrolador utilizado en AED. Se destaca que lo entregado por los pines digitales no debe superar los 5V y 20 mA.

Es importante especificar los componentes necesarios (ej: relés, TRIAC’s, relé de estado sólido, etc) y el diseño eléctrico del sistema. Investigue las diferentes alternativas, y justifique se elección (ventajas y desventajas)

\item \textit{Implementación y Programación:} Deben implementar el sistema y programar el microcontrolador para que regule las cargas eléctricas. Esto incluye la lectura de la temperatura mediante una RTD acondicionada. 

A su vez, debe setear con la computadora los máximos y mínimos de la temperatura, y conectar microcontrolador a tiempo real. Si ambos limites son cruzados, se deben activar dos alarmas específicas. Toda la información del sistema se tiene que visualizar en el display LCD.
\item \textit{Testeo:} Realizar pruebas exhaustivas del sistema en un entorno controlado para evaluar su capacidad para enfriar y calentar de manera efectiva. Deben documentar los resultados y analizar el rendimiento del sistema.

\item \textit{Entrega final:} Deben presentar un informe final que incluya la descripción detallada del diseño, implementación, esquemas eléctricos, resultados de las pruebas y conclusiones. 
También se debe presentar todo el conjunto armado y funcionando: Arduino + adecuador + entrenador +  placa desarrollada + cargas. No deben encontrarse sueltas, ya que puede llevar a un riesgo debido al manejo de alta potencia.
La misma será presentada y defendida el día de la Expotecnica 2023. 



\end{enumerate}
\centering
{\includegraphics[width=0.7\textwidth]{esq.png}\par}


\end{document}
